% Part 0 — The Kernel. Program-wide front matter, shared by all four volumes.
% Canonical source: kernel/part0_kernel.tex; copied into each volume by `make sync-style`.
% No \ref anywhere: this file compiles inside four different documents, so every
% cross-volume pointer is textual ([DL 5], [NLP 9], [SR 4], [CML 2]).
\chapter*{Part 0: The Kernel}
\addcontentsline{toc}{chapter}{Part 0: The Kernel}

Four volumes, several hundred questions. Almost all of it is downstream of about twenty-six ideas, and this part is those ideas.

They are not a summary. A summary compresses what the chapters say; the kernel is the layer \emph{above} them --- the generators from which their content can be rebuilt. Someone who holds these can answer questions this program never asked, because the answer becomes a derivation rather than a retrieval. That is what a Staff loop tests: not whether you have read about grouped-query attention, but whether, handed an unfamiliar serving constraint, you reach for the byte count.

\begin{tldr}
Read this part first and reread it last. Each principle carries a name, a statement, the places in the program where it does real work, and what an interviewer learns from your handling of it. When a chapter's material feels arbitrary --- a formula with no visible parent, a design choice presented as simply what people do --- go back to the kernel principle it instantiates. Arbitrariness is nearly always a missing generator, not a missing fact.
\end{tldr}

\section*{How to Use This Part}

\textbf{As a map.} The pointers under each principle run both ways: from here into the chapters for depth, and from a chapter back to here for the reason it exists.

\textbf{As a rehearsal script.} The unit of practice is the statement, spoken from a blank page in under a minute, plus two instantiations from different volumes.

\textbf{As a recovery move.} What costs most in a real interview is not ignorance but freezing on a forgotten constant. The kernel is what you say instead --- \emph{I do not recall the ratio, but the constraint is bandwidth, so what matters is bytes per token} --- and then rebuild. Interviewers reward this visibly: it is the clearest evidence that the knowledge is generative.

\begin{mathresult}{Five Expressions the Program Leans On}
\begin{align*}
\text{Loss} \ &= H(p) + \mathrm{KL}(p \,\|\, q) = \mathbb{E}_{p}[-\log q]
&& \text{Logit grad} \ = p - y \\
\text{Compute} \ &\approx 6ND, \ \text{optimal at } D \approx 20N
&& \text{KV per seq} \ = 2 L\, n_{kv} d_h s \ \text{bytes} \\
\text{Clock} \ &= \max(\text{FLOPs}/\text{peak}, \ \text{bytes}/\text{BW})
\end{align*}
\end{mathresult}

\begin{warningbox}
The kernel is not a checklist to recite. Saying ``Goodhart's law applies here'' without naming the specific proxy, the specific gap, and the specific mitigation is worse than saying nothing: it signals vocabulary without mechanism, which is exactly what senior interviewers are calibrated to detect. Each principle earns its place only when discharged into a concrete number, mechanism, or trade-off.
\end{warningbox}

\section*{Group A --- Objectives: What a Loss Actually Is}

Four principles about the thing you minimize. Many questions that look architectural are objective questions in disguise.

\subsection*{A1. Read every loss as a likelihood}

\textbf{Statement.} Cross-entropy training is maximum likelihood, and the loss is the data's expected code length, $H(p) + \mathrm{KL}(p\,\|\,q)$: there is an irreducible floor, every delta is bits saved, and choosing a loss is choosing a noise model.

\textbf{Where it bites.}
\begin{itemize}[itemsep=1pt,topsep=2pt,leftmargin=1.3em]
\item Perplexity as the exponential of per-token loss, and why loss cannot reach zero: the entropy floor of the text [NLP~2], [DL~23].
\item The taxonomy tying each loss to its noise model --- squared error to Gaussian, absolute error to heavier tails [DL~3] --- and the MLE/MAP framing behind it [CML~4].
\item Why label smoothing and distillation work on the \emph{target distribution}, not the optimizer [DL~3], [DL~14].
\end{itemize}

\textbf{The tell.} With it, a candidate derives a loss for an unfamiliar problem from its noise model; without it, they recite a menu and pick what they know.

\subsection*{A2. Start the derivation at the logits}

\textbf{Statement.} The gradient of softmax cross-entropy at the logits is exactly $p - y$: bounded, linear in miscalibration, free of the saturation softmax alone would introduce. Hence the default pairing, and hence logit-level interventions work at all.

\textbf{Where it bites.}
\begin{itemize}[itemsep=1pt,topsep=2pt,leftmargin=1.3em]
\item The whiteboard derivation with the log-sum-exp stabilization that must accompany it in code [DL~1], and the from-scratch implementation [DL~28].
\item Why squared error on a classifier vanishes precisely when the model is confidently wrong [DL~3].
\item The $T^2$ factor in distillation and temperature scaling for calibration [DL~3], [DL~23]; and z-loss, sensible only once you see cross-entropy constrains logit \emph{differences} and lets their scale drift [DL~15].
\end{itemize}

\textbf{The tell.} The signal is not the algebra but whether the result gets connected to numerical stability and calibration unprompted.

\subsection*{A3. Learn by comparison when labels run out}

\textbf{Statement.} InfoNCE is cross-entropy over a batch: identify the positive among the negatives. Temperature sets how sharply the hardest negatives dominate, and collapse --- all inputs to one point --- is the standing failure of any objective without them.

\textbf{Where it bites.}
\begin{itemize}[itemsep=1pt,topsep=2pt,leftmargin=1.3em]
\item SimCLR, MoCo and CLIP, and the collapse prevention in BYOL, SimSiam and DINO [DL~8]; spectral diagnostics for a suspect embedding space [DL~6].
\item Two-tower retrieval on in-batch sampled softmax, and why popularity-skewed negatives need the logQ correction [SR~4], [SR~6].
\item Hard-negative mining and its false-negative trap [DL~7]; word2vec negative sampling, the same objective a decade earlier [NLP~3].
\end{itemize}

\textbf{The tell.} Seen as classification-over-a-batch, batch size, temperature and logQ are three consequences of one equation; otherwise they are three unrelated tricks.

\subsection*{A4. Generate by learning a transport}

\textbf{Statement.} Every deep generative model maps a simple distribution onto the data and pays somewhere in the trilemma of quality, speed and tractable likelihood: autoregressive takes the exact factorization, VAEs a Jensen bound, diffusion a denoiser per noise level, GANs no likelihood.

\textbf{Where it bites.}
\begin{itemize}[itemsep=1pt,topsep=2pt,leftmargin=1.3em]
\item The ELBO and the reparameterization trick, and why diffusion displaced GANs: many small well-posed regressions instead of one minimax [DL~11].
\item Classifier-free guidance as an explicit diversity-for-fidelity trade, latent diffusion as the move that buys back the compute [DL~11], [DL~18].
\item The autoregressive branch is the language-model objective, so decoding strategy is a question about the same distribution [NLP~10].
\end{itemize}

\textbf{The tell.} ``Why diffusion and not a GAN?'' is answered well by reaching for training stability and mode coverage, badly by reaching for benchmark numbers.

\section*{Group B --- Scale: What Compute Buys}

Four principles about resources and capability --- what separates planning a run from describing one.

\subsection*{B1. Spend compute as 6ND}

\textbf{Statement.} Training costs about $6ND$ FLOPs, and at fixed budget the loss-minimizing allocation grows $N$ and $D$ together near $D \approx 20N$; inference-dominated economics break that ratio, because serving cost is multiplied by query volume.

\textbf{Where it bites.}
\begin{itemize}[itemsep=1pt,topsep=2pt,leftmargin=1.3em]
\item Budget allocation in both directions, and the MFU accounting that turns FLOPs into GPU-hours and dollars [DL~2], [DL~25], [DL~26].
\item Model sizing as a total-cost rather than training-cost decision, which is what moves the optimum [DL~24].
\item Scaling-law reasoning as asked in language-model rounds, including why early flagships were undertrained [NLP~7].
\end{itemize}

\textbf{The tell.} Ask for the FLOPs to train 70B on 1.4T tokens: with this principle the answer arrives in seconds, arithmetic exposed.

\subsection*{B2. Extrapolate the loss, not the benchmark}

\textbf{Statement.} Test loss falls as a power law in parameters, data and compute over many orders of magnitude --- how a large run is de-risked before it is paid for. But an exact-match metric over $k$ tokens goes as $p^k$, turning smooth progress into an apparent jump.

\textbf{Where it bites.}
\begin{itemize}[itemsep=1pt,topsep=2pt,leftmargin=1.3em]
\item The parametric forms, the IsoFLOP procedure, and what scales predictably and what does not [DL~2].
\item Emergence, and the case that much of it is an artifact continuous metrics smooth away [DL~2], [NLP~7].
\item Why a benchmark delta needs a significance argument before it means anything [DL~23], [NLP~10].
\end{itemize}

\textbf{The tell.} A weak answer to ``predict the final loss before spending the compute'' describes scaling laws; a strong one fits small runs and names what will not transfer.

\subsection*{B3. Distrust the bias--variance U-curve}

\textbf{Statement.} Past the interpolation threshold test error descends again: among the many zero-training-loss solutions, SGD's implicit bias selects low-norm, flat ones. The classical picture covers the underparameterized regime only.

\textbf{Where it bites.}
\begin{itemize}[itemsep=1pt,topsep=2pt,leftmargin=1.3em]
\item Double descent's three regimes, lottery tickets, grokking, and why larger models generalize better [DL~2].
\item The decomposition itself, still right for tabular problems with real train--test gaps [CML~4], and the cross-validation failure modes that decide whether you can measure it [CML~5].
\item Why an unconstrained GBDT overfits while a one-epoch pretraining run does not --- one word, two regimes [CML~2].
\end{itemize}

\textbf{The tell.} ``Your model hits 100\% training accuracy --- is it overfitting?'' separates answering from the regime you are in from answering from the textbook.

\subsection*{B4. Spend inductive bias only where data is scarce}

\textbf{Statement.} An architectural prior constrains the hypothesis space: the right constraint buys sample efficiency, and any constraint the data could have taught is wasted capacity at scale. The decision-relevant quantity is the crossover.

\textbf{Where it bites.}
\begin{itemize}[itemsep=1pt,topsep=2pt,leftmargin=1.3em]
\item Vision transformers against convnets and where the pretraining-scale crossover sits [DL~10]; why transformers displaced recurrent models on path length and parallelism [NLP~4].
\item The tabular question, where the tree prior still wins and the honest answer says why [CML~2], [DL~24].
\item Graph networks as a prior warranted only when the graph carries what a feature vector cannot [DL~13].
\end{itemize}

\textbf{The tell.} Asked to pick an architecture for a new domain, this principle makes you ask how much data exists before naming anything.

\section*{Group C --- Optimization: How the Descent Behaves}

Three principles about training dynamics: the difference between tuning by folklore and debugging by mechanism.

\subsection*{C1. Treat $\eta/B$ as the temperature}

\textbf{Statement.} Minibatch gradient noise has covariance proportional to learning rate over batch size and regularizes toward flat minima; the ratio, not either quantity alone, sets the dynamics --- hence the linear scaling rule, warmup, and a critical batch size.

\textbf{Where it bites.}
\begin{itemize}[itemsep=1pt,topsep=2pt,leftmargin=1.3em]
\item The linear scaling rule and its failure modes, and whether gradient accumulation truly equals a larger batch [DL~15].
\item Implicit regularization of SGD, and the flat-versus-sharp account of generalization [DL~2].
\item Diagnosing a run that spiked mid-training, where warmup, schedule and noise level are the first suspects [DL~25], [DL~15].
\end{itemize}

\textbf{The tell.} ``You doubled the batch size --- what else changes?'' is a one-line question whose good answer contains a ratio.

\subsection*{C2. See Adam as a diagonal preconditioner}

\textbf{Statement.} Adam divides each coordinate's step by an EMA of its gradient magnitude, approximating a diagonal Newton step; transformers need it because gradient scales differ by orders of magnitude across layers and embedding rows. AdamW pulls decay out of that rescaling, restoring its meaning as a Gaussian prior.

\textbf{Where it bites.}
\begin{itemize}[itemsep=1pt,topsep=2pt,leftmargin=1.3em]
\item The optimizer comparison and the decoupling argument, plus the eight bytes per parameter the two moments cost [DL~15].
\item Why layerwise learning-rate decay and width-scaling parametrizations exist: more attempts to equalize per-coordinate dynamics [DL~15], [DL~25].
\item $L_2$ regularization as a MAP prior --- the clean statement of what decoupling restores [CML~1], [CML~4].
\end{itemize}

\textbf{The tell.} Everyone states the update; the discriminating follow-up is what $v$ approximates and why that helps transformers specifically.

\subsection*{C3. Keep the residual highway clean}

\textbf{Statement.} Depth trains because the residual gives the gradient an identity path instead of a product of Jacobians, and normalization holds activations at a fixed scale; pre-norm leaves that path untouched, which is why modern language models are pre-norm with RMSNorm.

\textbf{Where it bites.}
\begin{itemize}[itemsep=1pt,topsep=2pt,leftmargin=1.3em]
\item Pre- versus post-norm, RMSNorm's dropped mean-centering, and why layer not batch normalization in a transformer [DL~4], [NLP~5].
\item Clipping, vanishing and exploding diagnosis, initialization scaling [DL~15]; the LSTM cell path as the same trick one generation earlier [NLP~4].
\item The residual stream as a shared bus --- the frame mechanistic interpretability works in [DL~20].
\end{itemize}

\textbf{The tell.} Saying residuals ``help gradients flow'' is the slogan; writing $1 + f'$ and explaining what the $1$ prevents is the principle.

\section*{Group D --- The Machine: Where the Time and the Bytes Go}

Six principles about hardware reality --- the largest group, because pricing things is where Staff candidates separate from strong senior ones.

\subsection*{D1. Ask compute-bound or memory-bound, always}

\textbf{Statement.} Time is the maximum of FLOPs over peak throughput and bytes over bandwidth, so a kernel's arithmetic intensity against the hardware ridge says which term you pay. Batch-one decode reads every weight per token and is bandwidth-stalled --- which is why nearly every inference optimization attacks bytes, not FLOPs.

\textbf{Where it bites.}
\begin{itemize}[itemsep=1pt,topsep=2pt,leftmargin=1.3em]
\item The roofline model, tiling and its failure modes, kernel fusion, and the profiling that decides which fix applies [DL~26].
\item FlashAttention as exact yet faster, never materializing the $n \times n$ matrix in HBM [DL~5], [DL~26].
\item Why quantization accelerates decode even at half precision, and why throughput rises almost free with batch size up to a point [DL~19].
\end{itemize}

\textbf{The tell.} Handed a performance question, this principle makes you ask which resource is saturated before proposing anything.

\subsection*{D2. Count bytes per parameter}

\textbf{Statement.} Mixed-precision Adam costs about sixteen to eighteen bytes per parameter before a single activation --- half-precision weights and gradients, an FP32 master copy, two moments --- with activations billed separately against batch, length, width and depth. Inference costs two bytes per parameter plus the cache.

\textbf{Where it bites.}
\begin{itemize}[itemsep=1pt,topsep=2pt,leftmargin=1.3em]
\item Why 7B exhausts an 80GB card in training but serves on a 24GB one; what each ZeRO stage shards; the worked 70B cluster layout [DL~15].
\item Which lines LoRA and QLoRA delete from the ledger --- the whole of why they are cheap [DL~16].
\item Gradient checkpointing as compute traded for memory, and why long context hurts training before serving [DL~15], [DL~25].
\end{itemize}

\textbf{The tell.} Memory questions are unusually honest signals: the arithmetic is short and cannot be bluffed --- you either enumerate the tensors or you do not.

\subsection*{D3. Do the KV-cache arithmetic}

\textbf{Statement.} The cache costs $2 \cdot L \cdot n_{kv} \cdot d_h \cdot s$ elements per sequence and is usually what binds serving batch size and context length --- not the weights. Once you can compute it, grouped-query attention, latent compression, paged allocation and quantized caches all read as one move at different points.

\textbf{Where it bites.}
\begin{itemize}[itemsep=1pt,topsep=2pt,leftmargin=1.3em]
\item The byte formula and the capacity plan it yields: how many concurrent requests fit on a node at a stated context [DL~19].
\item Why grouped-query attention barely moves quality but transforms serving economics, and what latent attention compresses instead [DL~5], [DL~14].
\item Prefix caching for shared system prompts, and why long-context serving stays expensive even with fast prefill [DL~19], [DL~21], [NLP~12].
\end{itemize}

\textbf{The tell.} The most common estimation question in a serving interview; the only acceptable answer is a number with its derivation.

\subsection*{D4. Separate prefill from decode}

\textbf{Statement.} One request holds two workloads: a compute-bound parallel prefill that sets time-to-first-token, and a memory-bound serial decode that sets tokens per second. Every serving technique targets one, and the SLO says which one you may spend on.

\textbf{Where it bites.}
\begin{itemize}[itemsep=1pt,topsep=2pt,leftmargin=1.3em]
\item Continuous batching against head-of-line blocking, chunked prefill and its interference cost, disaggregated fleets [DL~19], [DL~22].
\item Speculative decoding: why it preserves the target distribution exactly, and why it stops helping once the system is compute-bound [DL~19].
\item The latency budget of a retrieval-augmented or agentic turn, where retrieval, prefill and decode own differently priced slices [SR~8], [NLP~9], [DL~27].
\end{itemize}

\textbf{The tell.} ``First token is fast but generation is slow --- why?'' is one sentence with this principle, and unanswerable without it.

\subsection*{D5. Buy range before precision}

\textbf{Statement.} A float is sign plus exponent (range) plus mantissa (precision). Training breaks on range, not precision --- hence BF16, whose seven mantissa bits are swamped by gradient noise --- while inference tolerates aggressive weight quantization, because decode is bandwidth-bound and fewer bytes is proportionally faster.

\textbf{Where it bites.}
\begin{itemize}[itemsep=1pt,topsep=2pt,leftmargin=1.3em]
\item Loss scaling for FP16, why the FP32 master copy is kept, what breaks first at FP8 [DL~15], [DL~26].
\item Post-training against quantization-aware training, group-wise schemes and their overhead, and where quality falls off [DL~14], [DL~19].
\item Product quantization in an ANN index --- the same bytes-for-recall trade one layer down the stack [SR~3].
\end{itemize}

\textbf{The tell.} ``FP16 or BF16?'' is a bit-layout question; the strong answer names the exponent, then says why the lost mantissa bits do not matter.

\subsection*{D6. Split the model three ways}

\textbf{Statement.} Data parallelism replicates and all-reduces gradients, tensor parallelism splits matmuls and must stay inside a node, pipeline parallelism splits layers cheaply but pays a bubble. Real runs compose all three with sharded state; mixture-of-experts adds a fourth axis decoupling parameters from active FLOPs.

\textbf{Where it bites.}
\begin{itemize}[itemsep=1pt,topsep=2pt,leftmargin=1.3em]
\item The ``lay out a 70B run on 256 GPUs'' question, and replicated versus sharded data parallelism: what is stored where, sent when [DL~15], [DL~25].
\item The communication arithmetic pinning tensor parallelism to the intra-node interconnect, and the collectives behind it [DL~26].
\item Expert routing, load-balancing losses, and why sparse models are bandwidth-hungry despite low active FLOPs [DL~14], [NLP~7].
\end{itemize}

\textbf{The tell.} Candidates who have read about parallelism list the three types; those who have used it say which they would drop first, and why.

\section*{Group E --- Architecture and Adaptation: What to Compute, Precompute, and Change}

Five principles about structural choices --- each a question of where computation is allowed to happen.

\subsection*{E1. Read attention as a soft lookup}

\textbf{Statement.} $\mathrm{softmax}(QK^\top/\sqrt{d})V$ is a differentiable dictionary lookup --- queries score keys, the output averages values by those scores --- buying content-based routing at constant path length and quadratic cost, with $\sqrt{d}$ keeping logits out of saturation.

\textbf{Where it bites.}
\begin{itemize}[itemsep=1pt,topsep=2pt,leftmargin=1.3em]
\item The derivation with the variance argument for $\sqrt{d}$ [NLP~5], and the quantitative view: parameter and FLOP counts, and the length-versus-width crossover where attention overtakes the feed-forward block [DL~5].
\item Why KV caching is legal at all --- causal masking freezes past keys and values [DL~19]; implementing the block from scratch, masks included [DL~28].
\item What linear attention and state-space models surrender without the softmax lookup: precise retrieval from long context [DL~14], [NLP~7].
\end{itemize}

\textbf{The tell.} Everyone draws the block; the separating question is what the softmax is \emph{for} --- a relaxed argmax over a dictionary.

\subsection*{E2. Inject position deliberately}

\textbf{Statement.} Attention is permutation-equivariant, so position is a modeling choice. Rotary embeddings rotate queries and keys by position-dependent angles so scores depend only on relative offset, and because that offset rides in a phase, rescaling frequencies is what makes context extensible.

\textbf{Where it bites.}
\begin{itemize}[itemsep=1pt,topsep=2pt,leftmargin=1.3em]
\item Sinusoidal, learned, relative, rotary and additive-bias schemes, and the extrapolation trade each makes [DL~5], [NLP~5].
\item Position interpolation and frequency-aware rescaling behind 4K-to-128K extension, and the lost-in-the-middle behavior that survives it [DL~21], [NLP~7].
\item The long-context-versus-retrieval decision, priced with the cache arithmetic from D3 [DL~21].
\end{itemize}

\textbf{The tell.} Asked to extend context, this principle produces rotation phases going out of distribution; without it, fine-tuning on longer documents and hoping.

\subsection*{E3. Blame the tokenizer before the model}

\textbf{Statement.} Subword tokenization is a compression scheme fitted to a corpus, and it fixes the model's atomic units --- so context cost, multilingual pricing, and a family of notorious failures (character counting, digit handling, code whitespace) are tokenizer artifacts, not evidence about reasoning.

\textbf{Where it bites.}
\begin{itemize}[itemsep=1pt,topsep=2pt,leftmargin=1.3em]
\item BPE against WordPiece and Unigram, fertility, and the vocabulary-size trade against the embedding and output matrices [NLP~12], [DL~9]; the from-scratch drill [DL~28].
\item Tokenizer decisions at the level of a pretraining run, and their interaction with embedding-table memory and undertrained rare rows [DL~25], [DL~6].
\item The fairness consequence: identical content costing several times more tokens in a low-resource language [NLP~12].
\end{itemize}

\textbf{The tell.} The strawberry question is a gift: it reveals instantly whether someone reasons about the input representation or anthropomorphizes the model.

\subsection*{E4. Precompute what factorizes; rerank the rest}

\textbf{Statement.} If a score factorizes into a dot product of independently computed embeddings, the corpus can be precomputed and searched with an index; if not, the model runs once per pair. That distinction generates the funnel: cheap decomposable retrieval over millions, expensive joint scoring over the surviving dozens.

\textbf{Where it bites.}
\begin{itemize}[itemsep=1pt,topsep=2pt,leftmargin=1.3em]
\item Two-tower against cross-encoder cost analysis, with poly-encoders and late interaction as the middle ground [DL~7], [SR~4].
\item Index internals and the recall--latency--memory triangle that sets the retrieval operating point [SR~3].
\item The ads funnel from candidate generation through pre-rank, rank and auction, with a latency budget per stage [SR~6], [SR~8].
\item RAG as the same skeleton with a generator on the end --- which is why its failures are usually retrieval failures [NLP~9], [DL~21].
\end{itemize}

\textbf{The tell.} This turns the modal system-design question into a derivation: draw the funnel, then defend each stage's candidate count arithmetically.

\subsection*{E5. Adapt in a low-rank subspace}

\textbf{Statement.} The weight change a downstream task needs has low effective rank, so $\Delta W = BA$ with $r \ll d$ trains a fraction of a percent of parameters at near-full quality --- and the frozen base takes its gradients and moments off the memory ledger, which is the real saving.

\textbf{Where it bites.}
\begin{itemize}[itemsep=1pt,topsep=2pt,leftmargin=1.3em]
\item The low-rank hypothesis and its evidence, rank and target-module choices, and when full fine-tuning wins [DL~16], [NLP~6].
\item Quantized base weights with half-precision adapters, and multi-tenant serving off one base [DL~16], [DL~19].
\item The prompt, retrieve, or fine-tune decision, where this supplies the cost side [DL~24], [NLP~9].
\end{itemize}

\textbf{The tell.} ``Why does LoRA work?'' has a hypothesis for an answer, not a mechanism; saying so, and naming the evidence, is rarer than reciting the equation.

\section*{Group F --- Proxies: Why Deployed Systems Fail}

Four principles about the gap between what you optimize and what you want. They carry disproportionate weight at Staff level, because this is the part of the job that never appears in a paper.

\subsection*{F1. Assume the proxy will be gamed}

\textbf{Statement.} Every objective you can optimize is a proxy for something you cannot measure, and pressure lands preferentially where proxy and goal diverge. Mitigations share one shape: bound the pressure, keep a measure you never optimize against, iterate the proxy.

\textbf{Where it bites.}
\begin{itemize}[itemsep=1pt,topsep=2pt,leftmargin=1.3em]
\item Reward hacking after preference training --- sycophancy, verbosity, format exploitation --- and why it is structural, not a bug awaiting a patch [DL~17], [NLP~8].
\item Engagement optimization in ranked feeds, where the proxy drifts from satisfaction and the feedback loop amplifies the drift [SR~6], [DL~12].
\item Benchmark contamination and decontamination as pretraining-data discipline [DL~25]; judge drift, once the evaluator is itself the target [NLP~10], [DL~23].
\end{itemize}

\textbf{The tell.} The strongest answer concedes reward hacking cannot be permanently solved, then describes containment; the weakest proposes a better reward model.

\subsection*{F2. Bound the optimization pressure}

\textbf{Statement.} Preference optimization maximizes a learned reward minus a KL penalty against the reference policy. That penalty is the mechanism, not a nicety: it keeps the policy inside the distribution the reward model was fit on. Its constrained optimum has a closed form, and inverting it gives the direct-preference objective.

\textbf{Where it bites.}
\begin{itemize}[itemsep=1pt,topsep=2pt,leftmargin=1.3em]
\item The pipeline with its per-token penalty, and deriving DPO from the same objective through the Bradley--Terry likelihood [DL~17], [NLP~8].
\item What the coefficient trades: alignment tax against reward gain, and the diversity collapse from setting it too low [DL~17], [NLP~8].
\item Group-relative policy optimization and verifiable rewards, where a checkable reward shrinks the exploitable surface without closing it [DL~17], [NLP~8].
\end{itemize}

\textbf{The tell.} ``What happens without the KL term?'' is the cleanest probe of whether someone understands alignment training or has only seen its pipeline diagram.

\subsection*{F3. Choose the metric as a modeling decision}

\textbf{Statement.} A metric is a lossy projection of behavior onto a scalar, and picking it is itself a modeling decision: threshold-free measures flatter under heavy imbalance, ranking needs rank-aware metrics, downstream decisions need calibrated probabilities rather than correct orderings, and every comparison needs a significance argument.

\textbf{Where it bites.}
\begin{itemize}[itemsep=1pt,topsep=2pt,leftmargin=1.3em]
\item ROC against precision--recall under imbalance, calibration and temperature scaling, significance and multiple comparisons [DL~23], [CML~4].
\item The discount inside NDCG, the judgment supply chain feeding it, and interleaving as a ranking-specific design [SR~7].
\item Power, sample size, variance reduction, and the cost of peeking at a running test [CML~6].
\item Click-through calibration in an ad auction, where money depends on the probability, not the rank [SR~6].
\end{itemize}

\textbf{The tell.} ``The model is 99\% accurate'' should trigger an immediate question about base rates; if it does not, modeling depth cannot compensate.

\subsection*{F4. Assume training and serving disagree}

\textbf{Statement.} Models fail in production not because they are wrong but because the serving world differs: features computed by different code, joins that leaked information unavailable at prediction time, labels logged from what the old model chose to show, distributions that move after launch. Teacher forcing against free-running generation is the same mismatch, inside the model.

\textbf{Where it bites.}
\begin{itemize}[itemsep=1pt,topsep=2pt,leftmargin=1.3em]
\item The skew taxonomy, feature stores, point-in-time-correct joins, drift detection, retraining triggers [DL~22], [NLP~12].
\item Position bias and the logging feedback loop, with unbiased learning-to-rank and propensity weighting as the repair [SR~5], [SR~7].
\item The leakage taxonomy, and validation setups that hide leakage rather than reveal it [CML~5], [CML~6].
\item Exposure bias: teacher-forced training against autoregressive serving, with hallucination as one symptom [NLP~4], [NLP~13].
\end{itemize}

\textbf{The tell.} ``Offline metrics improved and online metrics dropped'' is the canonical Staff production question; a good answer enumerates causes by prior probability rather than guessing one.

\section*{The Kernel on One Page}

Reread this the night before. Each row is a name and the half-line that should recover the principle.

{\footnotesize
\begin{tabularx}{\textwidth}{@{}p{0.24\textwidth} X@{}}
\toprule
\textbf{Principle} & \textbf{Statement} \\
\midrule
\multicolumn{2}{@{}l}{\textit{A --- Objectives}} \\
Read loss as likelihood & Cross-entropy is MLE; the loss is code length; your loss is a noise model. \\
Start at the logits & Softmax cross-entropy gives $p - y$: bounded, non-saturating, calibration-shaped. \\
Learn by comparison & Cross-entropy over a batch; temperature weights hard negatives; collapse looms. \\
Generate by transport & One map from noise to data; a trilemma of quality, speed, tractable likelihood. \\
\addlinespace
\multicolumn{2}{@{}l}{\textit{B --- Scale}} \\
Spend compute as $6ND$ & Compute-optimal near $D \approx 20N$; serving volume pushes you past it. \\
Extrapolate the loss & Loss is a power law; thresholded metrics jump because $p^k$ is sharp in $p$. \\
Distrust the U-curve & Past interpolation, more capacity finds simpler solutions, not worse ones. \\
Bias where data is scarce & A prior is data you did not collect, and a ceiling you will eventually hit. \\
\addlinespace
\multicolumn{2}{@{}l}{\textit{C --- Optimization}} \\
$\eta/B$ is the temperature & Noise scales as learning rate over batch; the ratio sets the dynamics. \\
Adam preconditions & $m/\sqrt{v}$ equalizes per-coordinate steps; AdamW lifts decay out of it. \\
Keep the highway clean & Identity path plus fixed activation scale is why depth trains; pre-norm keeps it. \\
\addlinespace
\multicolumn{2}{@{}l}{\textit{D --- The machine}} \\
Compute or memory bound? & Compare arithmetic intensity to the hardware ridge before proposing a fix. \\
Count bytes per parameter & About 16--18 bytes to train, 2 to serve; activations are a separate bill. \\
Do the KV arithmetic & $2 L\, n_{kv} d_h s$ elements per sequence; this, not weights, caps batch. \\
Prefill is not decode & Compute-bound first token, bandwidth-bound throughput; each fix targets one. \\
Buy range before precision & Exponent bits keep training alive; mantissa bits are sold for bandwidth. \\
Split the model three ways & Data replicates, tensor splits matmuls in-node, pipeline splits layers. \\
\addlinespace
\multicolumn{2}{@{}l}{\textit{E --- Architecture and adaptation}} \\
Attention is a soft lookup & Queries score keys, output averages values; $\sqrt{d}$ prevents saturation. \\
Inject position by choice & Equivariance makes position a choice; rotary phase makes context extensible. \\
Blame the tokenizer first & Merges fix the atomic unit --- and with it cost, fairness, and spelling. \\
Factorize, then rerank & Decomposable scores buy an index; the rest buy accuracy one pair at a time. \\
Adapt in low rank & $\Delta W = BA$; task deltas are low-rank, so the frozen base leaves the ledger. \\
\addlinespace
\multicolumn{2}{@{}l}{\textit{F --- Proxies}} \\
The proxy will be gamed & Pressure lands where proxy and goal diverge; bound it, hold something back. \\
Bound the pressure & The KL term is that bound; inverting its optimum is where DPO comes from. \\
The metric is a choice & Imbalance, ranks, calibration and significance change what a number means. \\
Train and serve disagree & Different code, leaky joins, biased logs, drift --- and teacher forcing. \\
\bottomrule
\end{tabularx}
}

\begin{keyinsight}[The Move That Wins Rounds]
When you cannot recall a fact, say which principle governs it and rebuild from there, out loud. ``I do not remember the exact group size, but the constraint is bytes read per token, so the overhead is one scale and zero-point per group'' is a stronger answer than the memorized number, because it demonstrates the thing the number was only evidence for.
\end{keyinsight}
